% Part: first-order-logic
% Chapter: models-theories
% Section: expressing-props-of-structures

\documentclass[../../../include/open-logic-section]{subfiles}

\begin{document}

\olfileid{fol}{mat}{exs}

\olsection{表达 \printtoken{P}{structure} 的性质}

\begin{explain}
表达关于函数和关系的条件——或者更一般地说，表达结构中的函数和关系满足某些条件——常常是有用且重要的。
例如，我们希望能有办法将那些谓词符号“捕捉”到我们所期望的“含义”的结构，与那些未能做到这一点的结构区分开来。
当然，我们完全可以自由指定我们“意图”考虑的结构，例如，我们可以规定谓词符号~$\le$ 的解释必须是一个序，
或者声明我们只对这样的~$\Lang L$ 解释感兴趣：其论域由集合组成，且 $\Obj \in$ 被解释为“属于”关系。
但我们能用该语言本身的语句做到这一点吗？换句话说，
我们能用~$\Struct M$ 的语言中的语句（或一组语句）来表达~$\Struct M$ 的哪些条件？
有些条件是我们无法表达的。
例如，不存在~$\Lang L_A$ 的语句，它仅在满足 $\Domain M = \Nat$ 的结构~$\Struct M$ 中为真。
我们无法表达“论域仅包含自然数”。
但是，或许存在一些我们可以表达的结构的“结构性性质”。
我们能用语句表达结构的哪些性质？或者，换一种说法，哪些结构组成的集合，
可以被描述为使某个语句（或语句集）为真的结构？
\end{explain}

\begin{defn}[集合的模型]
设 $\Gamma$ 为语言~$\Lang L$ 中的一个语句集。
如果对所有的 $!A \in \Gamma$ 都有 $\Sat{M}{!A}$，则称结构~$\Struct M$ \emph{是~$\Gamma$ 的模型}。
\end{defn}

\begin{ex}
语句 $\lforall[x][x \le x]$ 在~$\Struct M$ 中为真，当且仅当 $\Assign{\le}{M}$ 是一个自反关系。
语句 $\lforall[x][\lforall[y][((x \le y \land y \le x) \lif x = y)]]$ 在~$\Struct M$ 中为真，当且仅当 $\Assign{\le}{M}$ 是反对称的。
语句 $\lforall[x][\lforall[y][\lforall[z][((x \le y \land y \le z)
      \lif x \le z)]]]$ 在~$\Struct M$ 中为真，当且仅当 $\Assign{\le}{M}$ 是传递的。
因此，\begin{align*}
\{\quad &\lforall[x][x \le x], \\
   & \lforall[x][\lforall[y][((x \le y \land y \le
    x) \lif x = y)]], \\
   &\lforall[x][\lforall[y][\lforall[z][((x \le y
      \land y \le z) \lif x \le z)]]] \quad \}
\end{align*} 的模型正是那些~$\Assign{\le}{M}$ 为自反、反对称且传递（即偏序）的结构。
于是，我们可以把它们当作\emph{偏序的一阶理论}的公理。
\end{ex}

\end{document}